%!TeX root=../thesis.tex

\chapter{Results and Refinements}
\label{chap:results}
\label{ch:results}

This chapter presents the empirical findings of the study and the tool refinements
they motivated, treating both as two faces of the same research contribution. The
\emph{Results} sections (Sections~\ref{sec:results-rq1}--\ref{sec:results-rq3}) report
how teams used Hypo-Stage, the architectural uncertainties they managed, and the
benefits and challenges they perceived. The \emph{Refinements} sections
(Section~\ref{sec:refinements}) follow directly and present each change made to
Hypo-Stage or its usage guidelines, linked to the specific observation that motivated it
and to the research question it helps answer (primarily RQ4). This integrated structure reflects the
action-research cycle at the heart of the study design: observation drives refinement,
and refinement shapes subsequent observation.

\section{Empirical Results (RQ1)}
\label{sec:results-rq1}

% TODO: Report patterns of Hypo-Stage and ArchHypo use.

\section{Empirical Results (RQ2)}
\label{sec:results-rq2}

% TODO: Characterize architectural uncertainties managed via Hypo-Stage.

\section{Empirical Results (RQ3)}
\label{sec:results-rq3}

% TODO: Report perceived benefits and challenges.

\section{Refinements (RQ4)}
\label{sec:results-rq4}
\label{sec:refinements}

% TODO: Summarize refinements emerging from the study (links to Section~\ref{sec:refinements}).

% TODO: Detail each refinement, its motivation, and linkage to research questions.
